\documentclass{article}

%maths
\usepackage{mathtools}
\usepackage{amsmath}
\usepackage{amssymb}
\usepackage{amsfonts}

%tikzpicture
\usepackage{tikz}
\usepackage{scalerel}
\usepackage{pict2e}
\usepackage{tkz-euclide}
\usetikzlibrary{calc}
\usetikzlibrary{patterns,arrows.meta}
\usetikzlibrary{shadows}

%pgfplots
\usepackage{pgfplots}
\pgfplotsset{compat=newest}
\usepgfplotslibrary{statistics}
\usepgfplotslibrary{fillbetween}

%colours
\usepackage{xcolor}

\usepackage{nicefrac}

\begin{document}
\section{Radicals and Complex Numbers}

\textbf{Radical Expression} is the form $\sqrt[n]{a}$. The $\sqrt{}$ is a radical sign and the expression under it is the \textit{radicand}, while the $n$ is called the \textit{index}.
\begin{center}
	\fbox{If $x = \sqrt[n]{a}$, then $x^n = a$}
\end{center}

\subsection{Adding and Subtracting Radical Expressions}

Radical expressions are called \textbf{like radical expressions} if the indexes and radicands are the same. \\
\\
Example: $3\sqrt{7} + 5\sqrt{7} = 8\sqrt{7}$

\subsection{Multiplying Radical Expressions}
To multiply we simply distribute each one then add them.\\
\\
\subsection{Dividing Radical Expressions}
To divide we simply apply the sqrt to both top and bottom.\\
\\
\subsection{Rational Exponents}
if n is a natural number greater than 1, and b is any nonnegative real number, then:\\
\\
\begin{center}
	$b^{1/n}=\sqrt[n](b)$
\end{center}

\subsection{Complex Numbers}
Imaginary values and pure imaginary numbers. The expression $\sqrt{-1}$ has no real answer. The symbol for it is \textit{i} because its called an imaginary value. Any expression that is a product of a real number with i is called a pure imaginary number. \\
\\
\begin{center}
	$i = \sqrt{-1}$, $ i^2 = -1$
\end{center}
\end{document}
